\documentclass[11pt,a4paper]{article}

\newcommand{\doctitle}{Delta Engine (MKM-DE-001)}
\newcommand{\docsubtitle}{Wolfram Oracle Benchmark of FS01, the Risky Annuity and Mark-to-Market}
\newcommand{\docversion}{0.2}
\newcommand{\docdate}{26-October-2026}
\newcommand{\docauthor}{David K Kelly}

% MKM Research Labs — Shared Document Header
% Usage: % MKM Research Labs — Shared Document Header
% Usage: % MKM Research Labs — Shared Document Header
% Usage: \input{../shared/mkm_header}  (before \begin{document})
%
% Required commands before \input:
%   \newcommand{\doctitle}{...}
%   \newcommand{\docsubtitle}{...}
%   \newcommand{\docversion}{1.0}
%   \newcommand{\docdate}{13-February-2026}
%   \newcommand{\docauthor}{David K Kelly}

% ─── Packages ───────────────────────────────────────────────────────────────
\usepackage[utf8]{inputenc}
\usepackage{amsmath,amsfonts,amssymb}
\usepackage{listings}
\usepackage{color,xcolor}
\usepackage{graphicx}
\usepackage{hyperref}
\usepackage[toc,page]{appendix}
\usepackage{tabto}
\usepackage{float}
\usepackage{comment}
\usepackage[en-GB]{datetime2}
\usepackage{setspace}
\usepackage{subcaption}
\usepackage{booktabs}
\usepackage{longtable}
\usepackage{adjustbox}
\usepackage{enumitem}
\usepackage{fancyhdr}
\usepackage{lastpage}

% ─── Page geometry ──────────────────────────────────────────────────────────
\onehalfspacing
\usepackage{geometry}
\geometry{a4paper, total={170mm,247mm}, left=20mm, top=25mm, bottom=25mm}

% ─── Headers and footers ───────────────────────────────────────────────────
\pagestyle{fancy}
\fancyhf{}
\fancyhead[L]{\small\textcolor{gray}{MKM Research Labs}}
\fancyhead[R]{\small\textcolor{gray}{\doctitle{} --- v\docversion}}
\fancyfoot[C]{\small\textcolor{gray}{CONFIDENTIAL}}
\fancyfoot[R]{\small\textcolor{gray}{Page \thepage\ of \pageref{LastPage}}}
\renewcommand{\headrulewidth}{0.4pt}
\renewcommand{\footrulewidth}{0.4pt}

% ─── Colours ────────────────────────────────────────────────────────────────
\definecolor{mkmblue}{HTML}{1565C0}
\definecolor{codegreen}{rgb}{0,0.6,0}
\definecolor{codegray}{rgb}{0.5,0.5,0.5}
\definecolor{codepurple}{rgb}{0.58,0,0.82}
\definecolor{backcolour}{rgb}{0.95,0.95,0.92}

% ─── Hyperref setup ─────────────────────────────────────────────────────────
\hypersetup{
    colorlinks=true,
    linkcolor=mkmblue,
    citecolor=mkmblue,
    urlcolor=mkmblue
}

% ─── Code listings ──────────────────────────────────────────────────────────
\lstdefinestyle{mystyle}{
    backgroundcolor=\color{backcolour},
    commentstyle=\color{codegreen},
    keywordstyle=\color{magenta},
    numberstyle=\tiny\color{codegray},
    stringstyle=\color{codepurple},
    basicstyle=\ttfamily\footnotesize,
    breakatwhitespace=false,
    breaklines=true,
    captionpos=b,
    keepspaces=true,
    numbers=left,
    numbersep=5pt,
    showspaces=false,
    showstringspaces=false,
    showtabs=false,
    tabsize=2
}
\lstset{style=mystyle}

% ─── Section formatting ────────────────────────────────────────────────────
\setcounter{tocdepth}{3}
\setcounter{secnumdepth}{3}

% ─── Shared macros ──────────────────────────────────────────────────────────
\newcommand{\mkmtitlepage}{%
    \begin{titlepage}
    \begin{center}
    \vspace*{1cm}
    \begin{figure}[H]\centering
        \includegraphics[width=0.4\textwidth]{../shared/Logos/MKM_logo.png}
    \end{figure}
    \vspace{1.5cm}
    {\Huge\bfseries\textcolor{mkmblue}{\doctitle}\par}
    \vspace{0.5cm}
    {\Large\textcolor{gray}{\docsubtitle}\par}
    \vspace{2cm}
    {\large MKM Research Labs\par}
    \vspace{0.5cm}
    {\large Version \docversion{} --- \docdate\par}
    \vspace{0.5cm}
    {\large \docauthor\par}
    \vspace{2cm}
    \textcolor{red}{\textbf{CONFIDENTIAL --- PROPRIETARY}}
    \end{center}
    \end{titlepage}
}

\newcommand{\mkmlegalpage}{%
    \clearpage
    \thispagestyle{empty}
    \begin{center}
    \large\textbf{Legal Notice}
    \vspace{1em}

    \textbf{PROPRIETARY AND CONFIDENTIAL}
    \end{center}

    \noindent This document, including all algorithms, models, methodology, formulas,
    calculations, and intellectual concepts contained herein, constitutes the exclusive
    intellectual property and confidential trade secrets of MKM Research Labs (``MKM'').

    \vspace{0.5em}
    \noindent Any usage, reproduction, distribution, modification, or disclosure of this
    document or any portion thereof, without the express prior written authorisation from
    MKM Research Labs is strictly prohibited and constitutes an infringement of intellectual
    property rights.

    \vspace{0.5em}
    \noindent The document is provided on an ``as is'' basis. MKM Research Labs makes no
    representations or warranties of any kind, express or implied, regarding its accuracy,
    reliability, or suitability for any particular purpose.

    \vspace{1em}
    \begin{center}
    All rights reserved. \copyright\ 2021--2026 MKM Research Labs.

    \vspace{0.5em}
    Governed by the laws of the United Kingdom.
    \end{center}
    \clearpage
}

\newcommand{\mkmhistorypage}[1]{%
    \section{Document History}
    \begin{table}[ht]
        \centering
        \begin{tabular}{l|l|l|l}
            \textbf{Date} & \textbf{Version} & \textbf{Description} & \textbf{Author} \\
            \hline
            #1
        \end{tabular}
        \label{tab:revision_history}
    \end{table}
    \clearpage
}

% ─── Bank validation protocol section commands ──────────────────────────────
% Use these to ensure every model doc has the required sections.
% They are just styled \section calls — content is provided by each doc.

\newcommand{\mkmexecsummary}{\section{Executive Summary}}
\newcommand{\mkmpurpose}{\section{Model Purpose and Scope}}
\newcommand{\mkmframework}{\section{Mathematical Framework}}
\newcommand{\mkminputs}{\section{Input Parameters and Calibration}}
\newcommand{\mkmimplementation}{\section{Implementation Details}}
\newcommand{\mkmsensitivity}{\section{Sensitivity Analysis}}
\newcommand{\mkmvalidation}{\section{Validation and Backtesting}}
\newcommand{\mkmtestresults}{\subsection{Automated Test Results}\input{test_results}}
\newcommand{\mkmlimitations}{\section{Model Limitations and Known Weaknesses}}

\newcommand{\mkmgovernance}{\section{Governance Validation Questions}
\noindent The following nine questions are mandated by Chapter~5 of the
\emph{Model Risk Governance for Vendors} handbook and must be explicitly
addressed for every model in the inventory. Response status is tracked in the
Model Governance Dashboard (MRC portal). \par\medskip}

\newcommand{\mkmvqitem}[3]{\subsection{Q#1 --- #2}
\noindent\textit{``#3''}\par\medskip\noindent}

\newcommand{\mkmchangelog}{\section{Change History}}
  (before \begin{document})
%
% Required commands before \input:
%   \newcommand{\doctitle}{...}
%   \newcommand{\docsubtitle}{...}
%   \newcommand{\docversion}{1.0}
%   \newcommand{\docdate}{13-February-2026}
%   \newcommand{\docauthor}{David K Kelly}

% ─── Packages ───────────────────────────────────────────────────────────────
\usepackage[utf8]{inputenc}
\usepackage{amsmath,amsfonts,amssymb}
\usepackage{listings}
\usepackage{color,xcolor}
\usepackage{graphicx}
\usepackage{hyperref}
\usepackage[toc,page]{appendix}
\usepackage{tabto}
\usepackage{float}
\usepackage{comment}
\usepackage[en-GB]{datetime2}
\usepackage{setspace}
\usepackage{subcaption}
\usepackage{booktabs}
\usepackage{longtable}
\usepackage{adjustbox}
\usepackage{enumitem}
\usepackage{fancyhdr}
\usepackage{lastpage}

% ─── Page geometry ──────────────────────────────────────────────────────────
\onehalfspacing
\usepackage{geometry}
\geometry{a4paper, total={170mm,247mm}, left=20mm, top=25mm, bottom=25mm}

% ─── Headers and footers ───────────────────────────────────────────────────
\pagestyle{fancy}
\fancyhf{}
\fancyhead[L]{\small\textcolor{gray}{MKM Research Labs}}
\fancyhead[R]{\small\textcolor{gray}{\doctitle{} --- v\docversion}}
\fancyfoot[C]{\small\textcolor{gray}{CONFIDENTIAL}}
\fancyfoot[R]{\small\textcolor{gray}{Page \thepage\ of \pageref{LastPage}}}
\renewcommand{\headrulewidth}{0.4pt}
\renewcommand{\footrulewidth}{0.4pt}

% ─── Colours ────────────────────────────────────────────────────────────────
\definecolor{mkmblue}{HTML}{1565C0}
\definecolor{codegreen}{rgb}{0,0.6,0}
\definecolor{codegray}{rgb}{0.5,0.5,0.5}
\definecolor{codepurple}{rgb}{0.58,0,0.82}
\definecolor{backcolour}{rgb}{0.95,0.95,0.92}

% ─── Hyperref setup ─────────────────────────────────────────────────────────
\hypersetup{
    colorlinks=true,
    linkcolor=mkmblue,
    citecolor=mkmblue,
    urlcolor=mkmblue
}

% ─── Code listings ──────────────────────────────────────────────────────────
\lstdefinestyle{mystyle}{
    backgroundcolor=\color{backcolour},
    commentstyle=\color{codegreen},
    keywordstyle=\color{magenta},
    numberstyle=\tiny\color{codegray},
    stringstyle=\color{codepurple},
    basicstyle=\ttfamily\footnotesize,
    breakatwhitespace=false,
    breaklines=true,
    captionpos=b,
    keepspaces=true,
    numbers=left,
    numbersep=5pt,
    showspaces=false,
    showstringspaces=false,
    showtabs=false,
    tabsize=2
}
\lstset{style=mystyle}

% ─── Section formatting ────────────────────────────────────────────────────
\setcounter{tocdepth}{3}
\setcounter{secnumdepth}{3}

% ─── Shared macros ──────────────────────────────────────────────────────────
\newcommand{\mkmtitlepage}{%
    \begin{titlepage}
    \begin{center}
    \vspace*{1cm}
    \begin{figure}[H]\centering
        \includegraphics[width=0.4\textwidth]{../shared/Logos/MKM_logo.png}
    \end{figure}
    \vspace{1.5cm}
    {\Huge\bfseries\textcolor{mkmblue}{\doctitle}\par}
    \vspace{0.5cm}
    {\Large\textcolor{gray}{\docsubtitle}\par}
    \vspace{2cm}
    {\large MKM Research Labs\par}
    \vspace{0.5cm}
    {\large Version \docversion{} --- \docdate\par}
    \vspace{0.5cm}
    {\large \docauthor\par}
    \vspace{2cm}
    \textcolor{red}{\textbf{CONFIDENTIAL --- PROPRIETARY}}
    \end{center}
    \end{titlepage}
}

\newcommand{\mkmlegalpage}{%
    \clearpage
    \thispagestyle{empty}
    \begin{center}
    \large\textbf{Legal Notice}
    \vspace{1em}

    \textbf{PROPRIETARY AND CONFIDENTIAL}
    \end{center}

    \noindent This document, including all algorithms, models, methodology, formulas,
    calculations, and intellectual concepts contained herein, constitutes the exclusive
    intellectual property and confidential trade secrets of MKM Research Labs (``MKM'').

    \vspace{0.5em}
    \noindent Any usage, reproduction, distribution, modification, or disclosure of this
    document or any portion thereof, without the express prior written authorisation from
    MKM Research Labs is strictly prohibited and constitutes an infringement of intellectual
    property rights.

    \vspace{0.5em}
    \noindent The document is provided on an ``as is'' basis. MKM Research Labs makes no
    representations or warranties of any kind, express or implied, regarding its accuracy,
    reliability, or suitability for any particular purpose.

    \vspace{1em}
    \begin{center}
    All rights reserved. \copyright\ 2021--2026 MKM Research Labs.

    \vspace{0.5em}
    Governed by the laws of the United Kingdom.
    \end{center}
    \clearpage
}

\newcommand{\mkmhistorypage}[1]{%
    \section{Document History}
    \begin{table}[ht]
        \centering
        \begin{tabular}{l|l|l|l}
            \textbf{Date} & \textbf{Version} & \textbf{Description} & \textbf{Author} \\
            \hline
            #1
        \end{tabular}
        \label{tab:revision_history}
    \end{table}
    \clearpage
}

% ─── Bank validation protocol section commands ──────────────────────────────
% Use these to ensure every model doc has the required sections.
% They are just styled \section calls — content is provided by each doc.

\newcommand{\mkmexecsummary}{\section{Executive Summary}}
\newcommand{\mkmpurpose}{\section{Model Purpose and Scope}}
\newcommand{\mkmframework}{\section{Mathematical Framework}}
\newcommand{\mkminputs}{\section{Input Parameters and Calibration}}
\newcommand{\mkmimplementation}{\section{Implementation Details}}
\newcommand{\mkmsensitivity}{\section{Sensitivity Analysis}}
\newcommand{\mkmvalidation}{\section{Validation and Backtesting}}
\newcommand{\mkmtestresults}{\subsection{Automated Test Results}\input{test_results}}
\newcommand{\mkmlimitations}{\section{Model Limitations and Known Weaknesses}}

\newcommand{\mkmgovernance}{\section{Governance Validation Questions}
\noindent The following nine questions are mandated by Chapter~5 of the
\emph{Model Risk Governance for Vendors} handbook and must be explicitly
addressed for every model in the inventory. Response status is tracked in the
Model Governance Dashboard (MRC portal). \par\medskip}

\newcommand{\mkmvqitem}[3]{\subsection{Q#1 --- #2}
\noindent\textit{``#3''}\par\medskip\noindent}

\newcommand{\mkmchangelog}{\section{Change History}}
  (before \begin{document})
%
% Required commands before \input:
%   \newcommand{\doctitle}{...}
%   \newcommand{\docsubtitle}{...}
%   \newcommand{\docversion}{1.0}
%   \newcommand{\docdate}{13-February-2026}
%   \newcommand{\docauthor}{David K Kelly}

% ─── Packages ───────────────────────────────────────────────────────────────
\usepackage[utf8]{inputenc}
\usepackage{amsmath,amsfonts,amssymb}
\usepackage{listings}
\usepackage{color,xcolor}
\usepackage{graphicx}
\usepackage{hyperref}
\usepackage[toc,page]{appendix}
\usepackage{tabto}
\usepackage{float}
\usepackage{comment}
\usepackage[en-GB]{datetime2}
\usepackage{setspace}
\usepackage{subcaption}
\usepackage{booktabs}
\usepackage{longtable}
\usepackage{adjustbox}
\usepackage{enumitem}
\usepackage{fancyhdr}
\usepackage{lastpage}

% ─── Page geometry ──────────────────────────────────────────────────────────
\onehalfspacing
\usepackage{geometry}
\geometry{a4paper, total={170mm,247mm}, left=20mm, top=25mm, bottom=25mm}

% ─── Headers and footers ───────────────────────────────────────────────────
\pagestyle{fancy}
\fancyhf{}
\fancyhead[L]{\small\textcolor{gray}{MKM Research Labs}}
\fancyhead[R]{\small\textcolor{gray}{\doctitle{} --- v\docversion}}
\fancyfoot[C]{\small\textcolor{gray}{CONFIDENTIAL}}
\fancyfoot[R]{\small\textcolor{gray}{Page \thepage\ of \pageref{LastPage}}}
\renewcommand{\headrulewidth}{0.4pt}
\renewcommand{\footrulewidth}{0.4pt}

% ─── Colours ────────────────────────────────────────────────────────────────
\definecolor{mkmblue}{HTML}{1565C0}
\definecolor{codegreen}{rgb}{0,0.6,0}
\definecolor{codegray}{rgb}{0.5,0.5,0.5}
\definecolor{codepurple}{rgb}{0.58,0,0.82}
\definecolor{backcolour}{rgb}{0.95,0.95,0.92}

% ─── Hyperref setup ─────────────────────────────────────────────────────────
\hypersetup{
    colorlinks=true,
    linkcolor=mkmblue,
    citecolor=mkmblue,
    urlcolor=mkmblue
}

% ─── Code listings ──────────────────────────────────────────────────────────
\lstdefinestyle{mystyle}{
    backgroundcolor=\color{backcolour},
    commentstyle=\color{codegreen},
    keywordstyle=\color{magenta},
    numberstyle=\tiny\color{codegray},
    stringstyle=\color{codepurple},
    basicstyle=\ttfamily\footnotesize,
    breakatwhitespace=false,
    breaklines=true,
    captionpos=b,
    keepspaces=true,
    numbers=left,
    numbersep=5pt,
    showspaces=false,
    showstringspaces=false,
    showtabs=false,
    tabsize=2
}
\lstset{style=mystyle}

% ─── Section formatting ────────────────────────────────────────────────────
\setcounter{tocdepth}{3}
\setcounter{secnumdepth}{3}

% ─── Shared macros ──────────────────────────────────────────────────────────
\newcommand{\mkmtitlepage}{%
    \begin{titlepage}
    \begin{center}
    \vspace*{1cm}
    \begin{figure}[H]\centering
        \includegraphics[width=0.4\textwidth]{../shared/Logos/MKM_logo.png}
    \end{figure}
    \vspace{1.5cm}
    {\Huge\bfseries\textcolor{mkmblue}{\doctitle}\par}
    \vspace{0.5cm}
    {\Large\textcolor{gray}{\docsubtitle}\par}
    \vspace{2cm}
    {\large MKM Research Labs\par}
    \vspace{0.5cm}
    {\large Version \docversion{} --- \docdate\par}
    \vspace{0.5cm}
    {\large \docauthor\par}
    \vspace{2cm}
    \textcolor{red}{\textbf{CONFIDENTIAL --- PROPRIETARY}}
    \end{center}
    \end{titlepage}
}

\newcommand{\mkmlegalpage}{%
    \clearpage
    \thispagestyle{empty}
    \begin{center}
    \large\textbf{Legal Notice}
    \vspace{1em}

    \textbf{PROPRIETARY AND CONFIDENTIAL}
    \end{center}

    \noindent This document, including all algorithms, models, methodology, formulas,
    calculations, and intellectual concepts contained herein, constitutes the exclusive
    intellectual property and confidential trade secrets of MKM Research Labs (``MKM'').

    \vspace{0.5em}
    \noindent Any usage, reproduction, distribution, modification, or disclosure of this
    document or any portion thereof, without the express prior written authorisation from
    MKM Research Labs is strictly prohibited and constitutes an infringement of intellectual
    property rights.

    \vspace{0.5em}
    \noindent The document is provided on an ``as is'' basis. MKM Research Labs makes no
    representations or warranties of any kind, express or implied, regarding its accuracy,
    reliability, or suitability for any particular purpose.

    \vspace{1em}
    \begin{center}
    All rights reserved. \copyright\ 2021--2026 MKM Research Labs.

    \vspace{0.5em}
    Governed by the laws of the United Kingdom.
    \end{center}
    \clearpage
}

\newcommand{\mkmhistorypage}[1]{%
    \section{Document History}
    \begin{table}[ht]
        \centering
        \begin{tabular}{l|l|l|l}
            \textbf{Date} & \textbf{Version} & \textbf{Description} & \textbf{Author} \\
            \hline
            #1
        \end{tabular}
        \label{tab:revision_history}
    \end{table}
    \clearpage
}

% ─── Bank validation protocol section commands ──────────────────────────────
% Use these to ensure every model doc has the required sections.
% They are just styled \section calls — content is provided by each doc.

\newcommand{\mkmexecsummary}{\section{Executive Summary}}
\newcommand{\mkmpurpose}{\section{Model Purpose and Scope}}
\newcommand{\mkmframework}{\section{Mathematical Framework}}
\newcommand{\mkminputs}{\section{Input Parameters and Calibration}}
\newcommand{\mkmimplementation}{\section{Implementation Details}}
\newcommand{\mkmsensitivity}{\section{Sensitivity Analysis}}
\newcommand{\mkmvalidation}{\section{Validation and Backtesting}}
\newcommand{\mkmtestresults}{\subsection{Automated Test Results}\input{test_results}}
\newcommand{\mkmlimitations}{\section{Model Limitations and Known Weaknesses}}

\newcommand{\mkmgovernance}{\section{Governance Validation Questions}
\noindent The following nine questions are mandated by Chapter~5 of the
\emph{Model Risk Governance for Vendors} handbook and must be explicitly
addressed for every model in the inventory. Response status is tracked in the
Model Governance Dashboard (MRC portal). \par\medskip}

\newcommand{\mkmvqitem}[3]{\subsection{Q#1 --- #2}
\noindent\textit{``#3''}\par\medskip\noindent}

\newcommand{\mkmchangelog}{\section{Change History}}

% T1 so the underscores in code identifiers print as underscores.
\usepackage[T1]{fontenc}
\usepackage{lmodern}
% Generated by tools.wolfram_oracle -- do not edit.
\newcommand{\bmGenerated}{2026-10-26}
\newcommand{\bmCodeFiles}{\texttt{src/models/trading/delta\_engine/pricer.py},\\ \texttt{src/models/trading/delta\_engine/engine.py}}
\newcommand{\bmCodeCommit}{\texttt{81c145f8f330}}
\newcommand{\bmCodeDirty}{}
\newcommand{\bmNumerics}{numpy 1.26.4, scipy 1.15.1, mpmath 1.4.1}
\newcommand{\bmWolfram}{15.0.0 for Mac OS X ARM (64-bit) (May 26, 2026)}
\newcommand{\bmWolframVia}{ via live kernel session (wolframclient)}
\newcommand{\bmWolframConfirms}{Wolfram and the 50-digit reference agree to $4.9\times10^{-25}$, so the reference is confirmed by a second, independent implementation.}
\newcommand{\bmCases}{120}
\newcommand{\bmChecks}{8}
\newcommand{\bmOutcome}{all pass}
\newcommand{\bmAnnuity}{$5.8\times10^{-16}$}
\newcommand{\bmMtm}{0.00}
\newcommand{\bmMtmLow}{0.00}
\newcommand{\bmNotional}{10{,}000{,}000}
\newcommand{\bmTradeSpread}{150}
\newcommand{\bmFd}{$3.0\times10^{-4}$}
\newcommand{\bmLinear}{0.94}
\newcommand{\bmTenor}{4.53}
\newcommand{\bmYieldLow}{0.4}
\newcommand{\bmYieldHigh}{2.4}
\newcommand{\bmHand}{2.724084}
\newcommand{\bmHandStated}{2.724084}
\newcommand{\bmConvexity}{0.32}
\newcommand{\bmMove}{50}
\newcommand{\bmFloorH}{0.0001}
\newcommand{\bmFloorModel}{0.0078}
\newcommand{\bmFloorHazard}{1.0039}
\newcommand{\bmFloorGrad}{1.0039}
\newcommand{\bmClaimsMet}{7}
\newcommand{\bmClaimsTested}{7}
\newcommand{\bmYieldShift}{100}
\newcommand{\bmTableTenor}{5}
\newcommand{\bmRates}{0.0001--0.5}

% The shared history table sets every column to its natural width; wrap the description.
\renewcommand{\mkmhistorypage}[1]{%
    \section{Document History}
    \begin{table}[ht]
        \centering
        \begin{tabular}{l|l|p{9cm}|l}
            \textbf{Date} & \textbf{Version} & \textbf{Description} & \textbf{Author} \\
            \hline
            #1
        \end{tabular}
    \end{table}
}

\begin{document}
% Long code paths and identifiers cannot hyphenate; let the spacing stretch instead.
\sloppy

\mkmtitlepage
\mkmlegalpage

\tableofcontents
\clearpage

% =========================================================================
\mkmexecsummary
% =========================================================================

This document benchmarks the Delta Engine (MKM-DE-001), which values and risk-manages
the PRS book: for every trade it computes the risky annuity, the flood sensitivity
FS01 and the mark-to-market. It runs in production, behind the trading routes, the
blotter, the XVA service and the historical end-of-day series. The fair spread it
differences is PR-001's, benchmarked separately.

\paragraph{The code computes its documented formulae, to the precision it reports.}
Over \bmCases{} cases --- annual hazard rates \bmRates, tenors 1 to 10 years, a flat
3\% curve and the book's yield curve, notional GBP~\bmNotional{} --- the risky annuity
agrees with a 50-digit reference to \bmAnnuity, and the spread change, FS01, the
returned annuity, the hazard sensitivity and mark-to-market agree to within half a
unit of the rounding the engine applies. \bmWolframConfirms

\paragraph{What changed since v0.1.} Benchmark v0.1 tested the engine against model
document v1.0 and found two defects in the code and three errors in the document:
\begin{enumerate}[leftmargin=*]
\item \textbf{MTM used a rounded annuity}, rounded to six places for display, which
  moved MTM by up to GBP~3.40 on GBP~10m. The engine now values with the unrounded
  annuity; MTM now agrees to the penny.
\item \textbf{FS01 was silent on the spread floor.} At $h=\bmFloorH$ the fair spread
  sits on the 2\,bp floor and FS01 reflects a \bmFloorModel\,bp spread change ---
  the right P\&L sensitivity, since MTM does not move either, but blind to the flood
  risk. The engine now flags the floor and reports the hazard sensitivity beside
  FS01: \bmFloorHazard\,bp, against an exact \bmFloorGrad\,bp.
\item \textbf{The document} said FS01 scaled linearly with the hazard rate (it is
  nearly flat), put the yield sensitivity at 3--5\% (it is 0.4--2.4\%), and gave the
  worked annuity as 2.738 (it is 2.724084). Model document v1.1 corrects all three,
  states that the basis delta \emph{is} the gauge delta, and documents the floor.
\end{enumerate}

\paragraph{The model document's claims} (v1.1, Table~\ref{tab:bm-claims}): all
\bmClaimsMet{} of the \bmClaimsTested{} claims tested hold, each to the precision the
document quotes: FS01 at 500\,bp is \bmLinear{} times FS01 at 50\,bp; a 100\,bp
yield rise lowers FS01 by \bmYieldLow--\bmYieldHigh\%; a 5-year trade carries
\bmTenor{} times the FS01 of a 1-year; the worked annuity is \bmHand; and the 1\,bp
bump is within \bmFd\,bp of the exact gradient. A \bmMove\,bp move is
\bmConvexity\% from full revaluation.

\paragraph{Evidence.} \bmChecks{} checks, outcome: \bmOutcome{}
(Table~\ref{tab:bm-checks}).
\begin{itemize}[leftmargin=*,nosep]
\item Wolfram: \bmWolfram{}\bmWolframVia.
\item Code under test, at commit \bmCodeCommit{}\bmCodeDirty:\\ \bmCodeFiles.
\item Numerics: \bmNumerics.
\item Generated \bmGenerated{} by \texttt{python -m tools.wolfram\_oracle de001}.
\end{itemize}

% =========================================================================
\mkmpurpose
% =========================================================================

\subsection{What is benchmarked}
\texttt{compute\_risky\_annuity}, \texttt{compute\_gauge\_delta} (FS01, the floor
flag and the hazard sensitivity), \texttt{compute\_basis\_delta} and
\texttt{compute\_mark\_to\_market} in \texttt{src/models/trading/delta\_engine/pricer.py},
called as \texttt{DeltaEngine.enrich\_trade} calls them, with the rounding they apply,
against model document v1.1 (\texttt{delta\_engine\_v1.1.pdf}).

\subsection{What is not}
The fair spread itself (MKM-PR-001's benchmark); the market-state inputs the trader
maintains; the risk-grid aggregation, which is a sum. This is a benchmark
of the implementation (validation question Q8), not independent validation (Q9).

% =========================================================================
\mkmframework
% =========================================================================

With annual hazard rate $h$, $\lambda=-\ln(1-h)$, quarterly dates $t_i=i/4$ and
discount $D(t)=e^{-r(t)t}$, the document defines
\begin{equation}
A(h,T)=\sum_{i=1}^{4T}\tfrac14\,e^{-\lambda t_i}D(t_i),\qquad
\text{FS01}=\delta\,\frac{s(h+10^{-4})-s(h)}{10^4}\,A(h,T)\,N,\qquad
\text{MTM}=\delta\,\frac{s-s_0}{10^4}\,A\,N,
\end{equation}
with $s$ PR-001's fair spread, floored at 2\,bp, $\delta=\pm1$ for a payer or receiver,
and $N$ the notional.

\paragraph{Why FS01 is flat in $h$.} On a flat curve PR-001's spread has the closed form
$s=10^4(1-R)\,(e^{\lambda\Delta t}-1)\,e^{r\Delta t/2}/\Delta t$, so
\begin{equation}
\frac{\mathrm{d}s}{\mathrm{d}h}=10^4\,(1-R)\,\frac{e^{\lambda\Delta t}\,e^{r\Delta t/2}}{1-h}
\approx 10^4 .
\end{equation}
A 1\,bp bump in $h$ moves the spread by about 1\,bp at any level, and FS01 is about
$10^{-4}A\,N$: it follows the annuity, which falls slowly as $h$ rises.

% =========================================================================
\section{Benchmark Method}
% =========================================================================

The equations above are evaluated with \texttt{mpmath} at 50 significant digits, from
the documents rather than the code; the exact gradient $\mathrm{d}s/\mathrm{d}h$ is a
high-precision central difference, checked against the closed form above. Wolfram
evaluates the same sums exactly, with the book curve as a piecewise function,
differentiates the spread symbolically, and evaluates the bumped spread at the very
double the engine computes, $h+10^{-4}$ in floating point. The engine rounds what it
returns, so each rounded output is judged to half a unit in its last place. MTM uses
an inception spread of \bmTradeSpread\,bp.

% =========================================================================
\mkmvalidation
% =========================================================================

% Generated by tools.wolfram_oracle -- do not edit.
\begin{table}[H]\centering\small
\caption{Benchmark checks}\label{tab:bm-checks}
\begin{adjustbox}{max width=\textwidth}
\begin{tabular}{p{82mm}p{24mm}rl}
\toprule
Check & Result & Tolerance & Outcome \\
\midrule
Risky annuity, unrounded (s3.1) & $5.8\times10^{-16}$ & $1.0\times10^{-12}$ & pass \\
Spread change $\Delta s$ (bp, to $10^{-4}$): model vs reference & $5.0\times10^{-5}$ & $5.0\times10^{-5}$ & pass \\
FS01 (GBP, to the penny): model vs reference & $5.0\times10^{-3}$ & $5.0\times10^{-3}$ & pass \\
Risky annuity as returned (to $10^{-6}$): model vs reference & $5.0\times10^{-7}$ & $5.0\times10^{-7}$ & pass \\
Mark-to-market, unrounded annuity (GBP, to the penny): model vs reference & $5.0\times10^{-3}$ & $5.0\times10^{-3}$ & pass \\
Hazard sensitivity, unfloored spread (bp, to $10^{-4}$): model vs reference & $5.0\times10^{-5}$ & $5.0\times10^{-5}$ & pass \\
Wolfram agrees with the reference: annuity, $\Delta s$, gradient & $4.9\times10^{-25}$ & $1.0\times10^{-12}$ & pass \\
Wolfram agrees with the reference: claims (annuity, FS01) & $1.4\times10^{-31}$ & $1.0\times10^{-12}$ & pass \\
\bottomrule
\end{tabular}
\end{adjustbox}
\end{table}


\subsection{The model document's claims}
% Generated by tools.wolfram_oracle -- do not edit.
\begin{table}[H]\centering\small
\caption{The model document's claims, tested}\label{tab:bm-claims}
\begin{adjustbox}{max width=\textwidth}
\begin{tabular}{p{38mm}p{40mm}p{56mm}l}
\toprule
Claim & Document says & Benchmark finds & Outcome \\
\midrule
s6 item 5: 1\,bp bump error & $<0.01$\,bp vs analytical gradient & $3.0\times10^{-4}$\,bp at most & met \\
s6 item 1: FS01 against $h$ & ``scales approximately linearly'', 50--500\,bp & FS01(500)/FS01(50) $= 0.94$, not 10 & \textbf{not met} \\
s6 item 2: tenor & 5Y $\approx 4.5\times$ 1Y & $4.53\times$ & met \\
s6 item 3: yields $+100$\,bp & FS01 changes 3--5\% & FS01 falls 0.4--2.4\% & \textbf{not met} \\
s7.2 hand calculation & $A\approx2.738$ & $A=2.724084$ & \textbf{not met} \\
s8 item 1: 50\,bp move & FS01$\times\Delta h$ differs from revaluation & by 0.32\% & finding \\
Spread floor (2\,bp) & not stated & at $h=0.0001$: $\Delta s=0.0078$\,bp vs gradient 1.004\,bp & finding \\
s3.5 basis delta & $\approx$ gauge delta & identical in every case: the basis argument is unused & finding \\
\bottomrule
\end{tabular}
\end{adjustbox}
\end{table}


\subsection{FS01 across the grid}
% Generated by tools.wolfram_oracle -- do not edit.
\begin{table}[H]\centering\small
\caption{FS01 in GBP from the engine{,} flat 3\% curve{,} $R=0${,} notional 10{,}000{,}000. Down a column it barely moves: FS01 is set by the annuity, not the hazard rate}\label{tab:bm-fs01}
\begin{adjustbox}{max width=\textwidth}
\begin{tabular}{rrrrrrr}
\toprule
$h$ & 1Y & 2Y & 3Y & 4Y & 5Y & 10Y \\
\midrule
0.0001 & 8 & 15 & 22 & 29 & 36 & 67 \\
0.0005 & 986 & 1{,}941 & 2{,}869 & 3{,}768 & 4{,}640 & 8{,}624 \\
0.001 & 986 & 1{,}942 & 2{,}868 & 3{,}766 & 4{,}637 & 8{,}609 \\
0.005 & 988 & 1{,}943 & 2{,}864 & 3{,}754 & 4{,}613 & 8{,}485 \\
0.01 & 991 & 1{,}944 & 2{,}859 & 3{,}738 & 4{,}583 & 8{,}334 \\
0.02 & 998 & 1{,}947 & 2{,}849 & 3{,}707 & 4{,}524 & 8{,}043 \\
0.05 & 1{,}018 & 1{,}956 & 2{,}820 & 3{,}618 & 4{,}353 & 7{,}252 \\
0.1 & 1{,}053 & 1{,}973 & 2{,}776 & 3{,}477 & 4{,}090 & 6{,}169 \\
0.2 & 1{,}135 & 2{,}017 & 2{,}701 & 3{,}233 & 3{,}645 & 4{,}673 \\
0.5 & 1{,}551 & 2{,}303 & 2{,}668 & 2{,}845 & 2{,}931 & 3{,}010 \\
\bottomrule
\end{tabular}
\end{adjustbox}
\end{table}


% =========================================================================
\mkmlimitations
% =========================================================================

\begin{enumerate}[leftmargin=*]
\item \textbf{Benchmark, not validation.} The engine is shown to compute its
  documented formulae; whether a parallel 1\,bp bump in the annual rate is the right
  risk measure is a separate question.
\item \textbf{The floor.} The engine now flags it, but risk limits and the risk grid
  aggregate FS01 alone; positions on the floor contribute almost nothing to them.
  Whether limits should also read the hazard sensitivity is for the desk and the
  model owner.
\item \textbf{Parallel bump only.} FS01 bumps the annual rate in parallel across
  tenors; term-structure and cross-gauge sensitivities are outside this benchmark.
\end{enumerate}

\mkmhistorypage{26-Oct-2026 & 0.1 & First issue, against model document v1.0 & David K Kelly \\
26-Oct-2026 & 0.2 & After the MTM and floor fixes, against model document v1.1 & David K Kelly \\}

\end{document}
